% -------------------------------------------------------------------------
% WBS ID: RES-MOD-SYS
% Deliverable: System Boundary and Mathematical Model Architecture
% Status: Draft complete for regional 2D/local 3D separation
% Evidence status: Regional and local hydrodynamic model boundaries established
% Review status: Not reviewed
% Last updated: 2026-07-18
% Dependencies: RES-PHY-SYN, RES-DAT-GEO, RES-GEO-REP, RES-MOD-EQS, RES-MOD-MUL
% Downstream interfaces: RES-MOD-LRG, RES-MOD-MUL, RES-NUM-SPEC, RES-VER-FRM
% -------------------------------------------------------------------------

\section{RES-MOD-SYS --- System Boundary and Mathematical Model Architecture}
\label{sec:res-mod-sys}

\subsection{Purpose and Scope}
\label{sec:res-mod-sys-purpose}

% TODO: Define what this Deliverable covers and explicitly excludes.

This Deliverable defines the continuous system boundary for the
regional tsunami model and the local three-dimensional
tsunami--barrier interaction model. It assigns physical responsibility
to each model without prescribing software classes, file formats or
implementation APIs.

\subsection{Research Questions}
\label{sec:res-mod-sys-research-questions}

% TODO: State the questions that must be answered before the Deliverable
% can be accepted.

The local-model integration answers:

\begin{enumerate}
  \item Which physical region is owned by the local three-dimensional
    model rather than the regional depth-averaged model?
  \item Which boundaries define the local air--water--barrier domain?
  \item Which outputs justify activating the local model?
  \item Which coupling modes must remain available between the
    regional and local models?
\end{enumerate}

\subsection{Terminology and Definitions}
\label{sec:res-mod-sys-definitions}

% TODO: Define the technical terminology, symbols, quantities and
% classifications used in this Deliverable.

\subsection{Literature Evidence}
\label{sec:res-mod-sys-literature-evidence}

% TODO: Record evidence from peer-reviewed papers, authoritative datasets,
% standards, technical reports and primary documentation.
%
% Do not create a detached literature-summary list. Explain how each source
% supports, contradicts or limits a project decision.

\textcite{qin2018comparison} showed that a three-dimensional RANS
model is needed when local vertical velocity structure, nonhydrostatic
pressure and direct fluid-force extraction are project outputs. Source
role: `3D-RANSVOF`. Project conclusion: the local barrier domain cannot
be represented solely by the regional depth-averaged equations.
Limitation: the result supports model separation and output
requirements, not the final numerical implementation.

\textcite{takase2016hybrid} and \textcite{romano2025coupling} support
separating regional and high-fidelity local domains while exchanging
state information across a controlled interface. Source role:
`3D-COUPLING`. Project conclusion: the architecture shall retain
separable regional and local model boundaries with configurable
interface placement. Limitation: neither paper directly fixes the
finite-volume transfer operators used by this project.

\subsection{Governing Relations and Quantitative Definitions}
\label{sec:res-mod-sys-governing-relations}

% TODO: Record relevant equations, dimensionless groups, empirical models,
% extraction rules or quantitative definitions.
%
% Use align environments for displayed equations.
% Every displayed equation must have a unique \label{eq:...}.
% Do not place punctuation after displayed equations.

\subsection{System Boundary}
\label{sec:res-mod-sys-system-boundary}

The regional model boundary is the horizontal two-dimensional
depth-averaged tsunami domain. It owns regional propagation,
bathymetric transformation, run-up and inundation at the selected
case-study scale. The local three-dimensional model owns
structure-scale pressure, overtopping, local free-surface topology and
FSI where depth-averaged assumptions are insufficient.

The local hydrodynamic domain is:

\begin{align}
  \Omega_{\mathrm{F}}
  &=
  \Omega_{\mathrm{L}}
  \setminus
  \Omega_{\mathrm{B}}
  \label{eq:res-mod-sys-local-fluid-domain}
\end{align}

where $\Omega_{\mathrm{L}}$ is the selected local coastal volume and
$\Omega_{\mathrm{B}}$ is the fixed candidate barrier or defence
geometry. The baseline local model is activated for run-up,
overtopping, breaking, local separation, pressure, shear traction,
force, moment and comparative barrier-performance analysis.

The local boundary is partitioned as:

\begin{align}
  \partial\Omega_{\mathrm{F}}
  &=
  \Gamma_{\mathrm{in}}
  \cup
  \Gamma_{\mathrm{out}}
  \cup
  \Gamma_{\mathrm{side}}
  \cup
  \Gamma_{\mathrm{top}}
  \cup
  \Gamma_{\mathrm{T}}
  \cup
  \Gamma_{\mathrm{B}}
  \label{eq:res-mod-sys-local-boundary-partition}
\end{align}

Here $\Gamma_{\mathrm{in}}$ is the regional-to-local incident-wave
boundary, $\Gamma_{\mathrm{out}}$ is the downstream open boundary,
$\Gamma_{\mathrm{side}}$ represents artificial lateral ocean cuts,
$\Gamma_{\mathrm{top}}$ is the atmospheric boundary,
$\Gamma_{\mathrm{T}}$ is fixed terrain and $\Gamma_{\mathrm{B}}$ is
the fixed barrier surface. The side boundaries represent unmodelled
lateral ocean, not physical walls, unless a later site-specific defence
geometry explicitly requires a wall boundary.

\subsection{State Variables and Parameters}
\label{sec:res-mod-sys-state-variables-parameters}

% TODO: Complete this RES-MOD Domain-specific subsection.

\subsection{Continuous Governing Model}
\label{sec:res-mod-sys-continuous-governing-model}

% TODO: Complete this RES-MOD Domain-specific subsection.

\subsection{Source Terms and Closures}
\label{sec:res-mod-sys-source-terms-closures}

% TODO: Complete this RES-MOD Domain-specific subsection.

\subsection{Initial, Boundary and Interface Conditions}
\label{sec:res-mod-sys-initial-boundary-interface-conditions}

% TODO: Complete this RES-MOD Domain-specific subsection.

\subsection{Model Coupling Requirements}
\label{sec:res-mod-sys-model-coupling-requirements}

The regional model shall provide boundary histories and checkpoint data
for local 3D/interface modelling. The final 2D--3D coupling operators,
FSI feedback rules and local extraction region remain separate
workstreams.

The Kamaishi and Sendai corridor polygons delivered by `RES-DAT-SCN`
are data/model interface objects. They determine where terrain,
observations and regional replay histories are clipped and sampled, but
they do not change the governing equations. A corridor boundary shall
therefore be justified by event-source location, target coast, severe
impact observations, defence geometry, topobathymetry coverage and
numerical relaxation requirements rather than by arbitrary drawing.

The baseline coupling path is one-way replay:

\begin{align}
  \mathcal{M}_{\mathrm{regional}}
  &\rightarrow
  \mathcal{M}_{\mathrm{local}}
  \label{eq:res-mod-sys-one-way-replay}
\end{align}

where a saved regional state and interface history drive the local
three-dimensional model without feedback during candidate screening.
A simultaneously coupled mode shall also be retained where reflection,
blocking or other local feedback materially changes regional or
nearshore quantities of interest.

\subsection{Sufficient Conditions for Validity}
\label{sec:res-mod-sys-sufficient-conditions-validity}

% TODO: Complete this RES-MOD Domain-specific subsection.

The regional-to-local separation is sufficient when the coupling
interface lies outside the region where vertical acceleration,
breaking, local impact pressure or barrier-generated recirculation
dominates. It shall be checked by moving the interface and comparing
the engineering quantities of interest.

\subsection{Candidate Approaches}
\label{sec:res-mod-sys-candidate-approaches}

% TODO: Define the credible candidate approaches considered.

\subsection{Critical Comparison}
\label{sec:res-mod-sys-critical-comparison}

% TODO: Compare candidates using physically and project-relevant criteria.
% Do not select an approach solely on popularity or implementation ease.

\subsection{Assumptions and Applicability}
\label{sec:res-mod-sys-assumptions}

% TODO: State assumptions and the conditions under which they are valid.

\subsection{Limitations and Failure Conditions}
\label{sec:res-mod-sys-limitations}

% TODO: State where the evidence, model or selected approach ceases to be
% reliable or applicable.

\subsection{Project-Specific Conclusions}
\label{sec:res-mod-sys-conclusions}

% TODO: Convert the evidence into conclusions specific to the Studentship
% project. Do not merely restate literature findings.

The project uses a two-model hydrodynamic architecture: a regional
depth-averaged NLSWE model for event-scale propagation and a local
three-dimensional URANS--VOF model for barrier-scale hydrodynamics.
The split is based on physical validity and output requirements rather
than software convenience.

\subsection{Selected Approach and Rationale}
\label{sec:res-mod-sys-selected-approach}

% TODO: Record the selected approach only after sufficient evidence exists.
% State the alternatives rejected and the reasons for rejection.

Use a regional 2D NLSWE model for the large horizontal domain and defer
local three-dimensional/interface/FSI modelling to the local domain.
This preserves the distinction between a two-dimensional horizontal
solve and cross-sectional visualisations extracted from that solve.

The selected local system role is:

\begin{quote}
  local three-dimensional incompressible two-phase
  Navier--Stokes/VOF model for fixed-terrain, fixed-barrier
  tsunami--barrier interaction, driven by regional NLSWE interface
  histories and used to extract free-surface, overtopping, pressure,
  shear, force, moment and energy metrics
\end{quote}

\subsection{Unresolved Decisions}
\label{sec:res-mod-sys-unresolved-decisions}

% TODO: Record unresolved questions, required evidence and decision owners.

Open system-boundary decisions are the final case-study local-domain
extent, the exact interface location, the downstream buffer length, and
the activation gate for simultaneous coupling.

\subsection{Requirements and Handoffs}
\label{sec:res-mod-sys-requirements}

% TODO: State requirements passed to other Research Domains, Software,
% verification activities or later execution work.

`RES-MOD-IBC` shall define the continuous boundary-condition roles on
each local boundary. `RES-NUM-IBC` shall convert those roles into
discrete boundary-face states and flux treatments. `RES-VER-FRM` shall
test that the chosen boundary and coupling placement do not contaminate
the local quantities of interest.

`RES-DAT-CAS` and Software shall provide the dataset manifest,
corridor-polygon identifiers and replay-boundary histories consumed by
the mathematical interface. The mathematical model shall treat those as
inputs, not as proof that validation has been executed.

\subsection{Deliverable Completion Record}
\label{sec:res-mod-sys-completion-record}

% TODO: Record whether the Deliverable is:
%
% - Not started
% - In progress
% - Draft complete
% - Reviewed
% - Accepted
%
% Acceptance requires:
%
% - the research questions to be answered;
% - claims to be supported by appropriate evidence;
% - assumptions and limitations to be explicit;
% - the project-specific conclusion to be stated;
% - downstream requirements to be traceable;
% - unresolved decisions to be closed or formally deferred.
